Nuclei instance segmentation in H\&E histopathology is essential for computational pathology tasks such as tumor-infiltrating lymphocyte (TILs) scoring and cancer grading. However, densely overlapping nuclei and the reliance of prior methods on handcrafted post-hoc processing (NMS, watershed) pose geometric and computational challenges that hinder deployment on resource-constrained edge devices. Existing methods remain mask-based and depend on such post-hoc processing to separate touching instances, making them unsuitable for lightweight edge inference.

This research introduces RayCastED, a lightweight, fully convolutional, end-to-end contour-based instance segmentation model that predicts star-convex polygons directly from an anchor grid without any post-hoc processing. RayCastED adopts a YOLO26-derived detection backbone with dual-label assignment for NMS-free inference, inter-scale competition to suppress cross-scale duplicate detections, ResoConv wavelet-based downsampling for boundary preservation, and a split classification head to address foreground-background imbalance. A novel analytical ray distance loss recomputes ground-truth rays from the predicted centroid, coupling centroid accuracy with boundary supervision and providing scale-invariant gradients.

Evaluations on PanNuke under 3-fold cross-validation show that RayCastED achieves competitive panoptic quality against substantially heavier baselines while attaining the most consistent per-tissue performance, and demonstrates viable real-time inference on the Jetson Orin Nano embedded platform. These results establish RayCastED as a post-hoc-processing-free contour-based model that unifies boundary accuracy, lightweight inference, and edge deployability for nuclei analysis. Code and models are available at \url{https://github.com/DanishRitonga/RayCastED}.

\noindent\textbf{Keywords} : Nuclei Instance Segmentation, Computational Pathology, Star-convex Polygons, Lightweight Architecture, Edge Deployment
